\honest{Raised in Technophilia}{Eliezer Yudkowsky}{2008}

My father used to say that if the present system had been in place a hundred years ago, automobiles would have been outlawed to protect the saddle industry. At nine I read Jerry Pournelle's \textsc{A Step Farther Out}. It was a reply to Paul Ehrlich and the Club of Rome, who said in the 1960s and 1970s that the Earth was running out of resources and that mass famine was years away; to Jeremy Rifkin's so-called fourth law of thermodynamics; and to the people trying to regulate nuclear power into oblivion. \nb{Ehrlich did forecast famine within years: ``In the 1970s hundreds of millions of people will starve to death.'' The Club of Rome's \textsc{The Limits to Growth} put its limits ``sometime within the next one hundred years.''}

I grew up with a clear picture of a battle fought over and over since the Industrial Revolution, with historical evidence about how it came out. On one side were the scientists and engineers who had driven every rise in living standards since the Dark Ages, and whose work supported democracy, an educated populace, a middle class and the outlawing of slavery. On the other were those who once opposed vaccination, anaesthesia in childbirth, steam engines and heliocentrism: theologians calling for a return to a perfect age that never existed, old politicians set in their ways, special interests that stood to lose, and people afraid of what they did not understand. In the middle were the ``pretenders to Deep Wisdom,'' who said, ``in defiance of brute historical fact,'' that science of itself was ``neither good nor evil,'' and who set up solemn committees to display their caution, as if the truth were always a compromise and anyone could see that far ahead. Would humanity have done better with a public debate on the adoption of fire, and committees to oversee its use? \nb{A record of good uses is no evidence against a claim about what science is of itself.} So I was allergic to anyone who said that technology has risks as well as benefits. I presumed such a person wanted cheap applause or meant to regulate the technology into oblivion.

The allergy has a reason. Today Robin Hanson wrote about the slow American approval of drugs already approved abroad. A commenter noted that thalidomide was sold in 50 countries, but little was given out in the United States, so of 10,000 malformed children worldwide only 17 were American. But how many people died waiting for the drugs that did not go wrong? And that count leaves out the drugs never developed because approval is long and costly. According to the source I link, slow approval of drugs in the United States prevents 5,000 casualties a year and causes 20,000 to 120,000. \nb{The figures appear to be Dale Gieringer's, which are per decade. The ratio is the same.} So there is ``a historical record showing over-conservativeness'': many silent deaths from regulation outweigh a few visible deaths from its absence. If you are really in the middle, why not say that technology has benefits as well as risks? One regulator, one kind of product, and the list of old controversies: that is the record.

Looking back, I think my childhood view was not ``too wrong'' about who the Bad Guys were, though they are ordinary people whose worldview puts them in the right, not evil mutants. But it is much easier to say what not to do than to get it right. My mistake was to think that avoiding everything the Bad Guys did made me a Good Guy. Seeking a middle way ``is usually wrong''; the Right Way is not a compromise with anything. My evidence is the pretenders. They smiled down on technophiles and technophobes alike as immature, and because of their example any departure from charging straight ahead felt like joining them.

The first crack came, I think, in 1997 or 1998, in a debate about nanotechnology. The question was whether offense would be easier than defense. People said defense would be easy and talked of unassailable diamondoid walls, for programmable matter, when we cannot even secure computer networks where we see every one and zero. I said diamond does not stop a nuclear weapon; offense has beaten defense since 1945, and nanotechnology did not look likely to change that. By the end I had noticed, for the first time, that the survival of ``Earth-originating intelligent life'' was at risk. \nb{The Extropians mailing list archive has this debate, from August 1997, with the young Yudkowsky's words: ``I am darn well entitled to demand that your immune system stand up against nuclear weapons.''} Until then I had thought only individual lives were at stake. I had not rejected the larger possibility; I had never seen it, and once the topic came up I did. I do not remember how that trick worked.

It may sound as if I was an idiot, but the truth is scarier. I was a sharp Traditional Rationalist, far above average. I knew that hypotheses must be testable and that rationalization is not allowed, I could play Rationalist's Taboo, and I was obsessed with self-awareness; I knew no Bayes or Kahneman. ``So what?'' Nature does not grade on a curve, and one step away from the Way can repeal all other protections. Even ``a lot of rationality'' does not pass ``the astronomically high barrier'' for things to work. That is why many people think intelligence is not everything, or that rationalists do no better in real life. \nb{Yudkowsky drew the same lesson in ``My Wild and Reckless Youth.''}

Do not blame Pournelle, my father or science fiction for my mistake. Pournelle said often that once you leave Earth, if you are careless sealing your pressure suit just once, you die. But that happened to minor characters; the hero rarely died halfway through. I filtered the teachings through hope: rationality and hope on one side, ignorance and despair on the other. I was reluctant to learn to drive, since cars looked unsafe, but I treated the Future as immortal: ``Lives could be lost, but not the Future.''

When I saw that nanotechnology could cause extinction, I thought, explicitly, that I must have been too attached to its benefits and had flinched from the thought of human extinction. Then I did not stop and rethink the conclusions built on my old attitude. My mind found reasons to keep the old plans, and I changed as little as I could. I decided that we had to race to AI and get there before nanotechnology: the same plan as before, with a new reason. \nb{The record agrees. In 1996 Yudkowsky's slogan was ``Get to the Singularity as fast as possible!'' In 1999 the reason given for the hurry was that the longer it took, ``the higher the chance of humanity wiping itself out.''} I guess most people are like that, and Traditional Rationality did not change it. It took a stronger boot to the head before I fully saw my mistake. To be continued.
